\documentclass[10pt,a4paper]{amsart}
\usepackage[margin=19mm]{geometry}
\usepackage{amsmath,amssymb,mathtools}
\usepackage[scr=rsfs,cal=euler]{mathalfa}
\usepackage{xfrac}
\usepackage[unicode,hidelinks,bookmarks=false]{hyperref}
\newcommand{\ie}{i.e.\ }
\newcommand{\cf}{cf.\ }
\newcommand{\resp}{respectively\ }
\newcommand{\Subsection}{\subsection}
\providecommand{\nobreakdash}{\nobreak\hbox{-}}
\setlength{\emergencystretch}{2em}
\multlinegap0pt
\usepackage{bm}
\usepackage{bbm}
\usepackage{dsfont}
\usepackage{xspace}
\usepackage{cancel}
\usepackage{mathtools}
\usepackage{amsthm}
\makeatletter
\@ifundefined{theorem}{\newtheorem{theorem}{Theorem}}{}
\@ifundefined{lemma}{\newtheorem{lemma}[theorem]{Lemma}}{}
\makeatother
\newcommand{\QFR}{\mathrm{QFR}}
\title[Quantum Fractional Revival and Entanglement Entropy in Unita]{Quantum Fractional Revival and Entanglement Entropy in Unitary Cayley Graphs}
\author{Duaa Abdullah}
\date{Source: arXiv:2605.13645v1 (CC BY 4.0)}
\begin{document}
\maketitle

\noindent\emph{Source and licence.} Quantum Fractional Revival and Entanglement Entropy in Unitary Cayley Graphs. Version arXiv:2605.13645v1 (CC BY 4.0), distributed under the Creative Commons Attribution 4.0 International licence, as recorded in the arXiv licence metadata for this exact version and shown on the abstract page https://arxiv.org/abs/2605.13645v1. Full rights notice: licenses/arxiv-2605.13645-CC-BY-4.0.txt.\par\smallskip

\noindent\textit{Editorial scope.} Counted unit: the Theorem whose source label is \texttt{Thm2revivaltime} in the article (source0.tex lines 449--455, source label \texttt{Thm2revivaltime}), delivered together with the article's own complete original proof of it (source0.tex lines 456--514, one \texttt{proof} environment of 59 lines). No mathematics is written, repaired or re-derived: statement and proof are the article's own text line for line. Required context retained uncounted: lem01eigs2p (lines 382--396). Every definition, display and label of those lines is retained unchanged. Omitted from this package and not redistributed: 1--381, 397--448, 515--1148. Nothing inside a retained excerpt is omitted or altered: every retained body is delivered exactly as the article's lines are written, so no omission receipt of any kind accompanies this package. Citations: the retained excerpts carry no citation command at all, so no bibliography entry of the article is transcribed and no reference is redistributed. Reference adaptation: none was needed and none was made; every reference inside the delivered unit resolves inside it, so no unresolved reference or citation is produced. Printed slips kept literal: no printed word, symbol, hypothesis or number of the counted statement or of its proof was altered. Layout and class: the article's own class is \texttt{amsart} with packages including fontenc, geometry, graphicx, amssymb, stmaryrd, enumitem, tikz, subcaption, float, url, mathtools, amsmath, amsfonts, amsthm; the wrapper is the lane's pinned \texttt{amsart} editorial wrapper at 10pt on A4, which restates the theorem-like environments the retained text needs and adds the article's own macro definitions that the retained text needs (\texttt{\textbackslash QFR}), with extra wrapper packages (bm, bbm, dsfont, xspace, cancel, mathtools). The delivered numbering is the unit's own. Assets: no retained span contains a figure environment, a graphics-inclusion command, a PDF-page inclusion, a table or a TikZ picture; no image file is copied into this package, the package declares no asset dependency and no image file is redistributed with it. Licence: version arXiv:2605.13645v1 of this article is distributed under the Creative Commons Attribution 4.0 International licence, as recorded in the arXiv licence metadata for this exact version and shown on the abstract page https://arxiv.org/abs/2605.13645v1. A full rights notice is delivered at licenses/arxiv-2605.13645-CC-BY-4.0.txt. Characters: every retained excerpt is pure ASCII, so the delivered engine, XeLaTeX with its default Latin Modern text font, reports no missing character.
\begin{lemma}\label{lem01eigs2p}
Let $p$ be an odd prime where $\varphi(n)=p-1$. Then, the adjacency eigenvalues
  of $X = (V(\mathbb{Z}_n),E(S))$ satisfying $ \lambda_\kappa=r(\kappa, 2p)$ where 
  \[
    \lambda_\kappa=
    \begin{cases}
      p-1,  & \text{if } p \mid \kappa \text{ and } 2 \nmid \kappa \text{ (equivalently, } \kappa = p), \\
      -(p-1), & \text{if } 2p \mid \kappa, \quad  \kappa= 0), \\
      1,   & \text{if } \gcd(\kappa, 2p) = 2, \quad 2 \mid \kappa,\; p \nmid \kappa), \\
      -1,  & \text{if } \gcd(\kappa, 2p) = 1, \quad \gcd(\kappa, 2p) = 1).
    \end{cases}
  \]
  More precisely, $\lambda_0 = p-1$, $\lambda_p = -(p-1)$, and for $\kappa \notin \{0, p\}$:
  $\lambda_\kappa = 1$ if $2 \mid \kappa$, and $\lambda_\kappa = -1$ if $2 \nmid \kappa$.
\end{lemma}

\begin{theorem}~\label{Thm2revivaltime}
Let $X = (V(\mathbb{Z}_{2p}), E(S))$ the unitary Cayley graph with an odd prime $n=2p$. Then $X$ admits $\QFR$ between pair $(0, p)$ under the adjacency Hamiltonian at the
  minimum time
\begin{equation}~\label{eqq1Thm2revivaltime}
 t^* = \frac{2\pi}{p}.
\end{equation}
\end{theorem}

\begin{proof}
Assume $p$ be an odd prime where $n=2p$. Then, $\QFR$ condition at time $t$ satisfied with 
  \begin{equation}~\label{eqq2Thm2revivaltime}
  \begin{aligned}
&\alpha = \frac{1}{n}\sum_{d=0}^{n-1} e^{i\lambda_{\kappa} t} \in \mathbb{C},\\
& \beta  = \frac{1}{n}\sum_{d=0}^{n-1} e^{i\lambda_{\kappa} t}\omega_n^{pd} \in \mathbb{C}\setminus\{0\},
  \end{aligned}
  \end{equation}
  where $|\alpha|^2 + |\beta|^2 = 1.$ Hence, according to Lemma~\ref{lem01eigs2p} and by considering that $\omega_n^{p\kappa}=(-1)^{\kappa}$, we established  that the eigenvalue type with partition $\{0, \ldots, n-1\}$ as
  \begin{equation}~\label{eqq3Thm2revivaltime}
   I_0 = \{0\},\; I_p = \{p\},\; I_+ = \{d : \gcd(d,2p)=1\},\; I_- = \{d : \gcd(d,2p)=2\},      
  \end{equation}
  where the terms holds in case $|I_+|=|I_-|=p-1$. Thus, we noticed that at $t = 2\pi/p$ satisfying 
  \begin{itemize}
    \item $e^{i\lambda_0 t} = e^{i(p-1)\cdot 2\pi/p}$.
    \item $e^{i\lambda_p t} = e^{-i(p-1)\cdot 2\pi/p}$.
    \item $e^{i\cdot 1 \cdot t} = e^{2\pi i/p}$ for each $d \in I_+$.
    \item $e^{i\cdot(-1)\cdot t} = e^{-2\pi i/p}$ for each $d \in I_-$.
  \end{itemize}
Therefore, 

\textbf{Case 1.} If $\alpha = \cos(2\pi/p) \in \mathbb{R}$. Then
  \begin{align*}
    2p\,\alpha &= e^{2\pi i(p-1)/p} + e^{-2\pi i(p-1)/p}
                 + (p-1)e^{2\pi i/p} + (p-1)e^{-2\pi i/p} \notag\\
               &= 2\cos\!\left(\frac{2\pi(p-1)}{p}\right) + 2(p-1)\cos\!\left(\frac{2\pi}{p}\right)\notag\\
               &= 2\cos\!\left(2\pi - \frac{2\pi}{p}\right) + 2(p-1)\cos\!\left(\frac{2\pi}{p}\right)\notag\\
               &= 2\cos\!\left(\frac{2\pi}{p}\right) + 2(p-1)\cos\!\left(\frac{2\pi}{p}\right)
                = 2p\cos\!\left(\frac{2\pi}{p}\right).
  \end{align*}
  Hence $\alpha = \cos(2\pi/p) \in \mathbb{R}$. 

  \medskip

\textbf{Case 2.} If $\beta= -i\sin(2\pi/p)$. Then, 
 for $\beta$, we find that $(-1)^{\kappa}=-1$ for
  \begin{itemize}
      \item  odd $\kappa$ contributing $I_+$ when $p$ is odd prime,
      \item and $\kappa$ odd non-multiple of $p$.
  \end{itemize}
Also, $(-1)^{\kappa}=1$ for even $\kappa$ contributing $I_-$. Thus
  \begin{align*}
    2p\,\beta &= (-1)^0 e^{2\pi i(p-1)/p} + (-1)^p e^{-2\pi i(p-1)/p}\\
&+ \sum_{d\in I_+}(-1)^d e^{2\pi i/p}  + \sum_{d\in I_-}(-1)^d e^{-2\pi i/p}.
  \end{align*}
Similarly, since $p$ is odd, $(-1)^p = -1$. It is elements of $I_+$ where odd $\kappa$ implies that $(-1)^{\kappa}=-1$; and it is elements
  of $I_-$ where even $\kappa$ implies that $(-1)^{\kappa}=1$. Thus:
  \begin{align*}
    2p\,\beta &= e^{2\pi i(p-1)/p} - e^{-2\pi i(p-1)/p}
                 - (p-1)e^{2\pi i/p} + (p-1)e^{-2\pi i/p}\notag\\
               &= 2i\sin\!\left(\frac{2\pi(p-1)}{p}\right) - 2i(p-1)\sin\!\left(\frac{2\pi}{p}\right)\notag\\
               &= -2i\sin\!\left(\frac{2\pi}{p}\right) - 2i(p-1)\sin\!\left(\frac{2\pi}{p}\right)
                = -2ip\sin\!\left(\frac{2\pi}{p}\right).\label{eq:beta_calc}
  \end{align*}
  Hence $\beta= -i\sin(2\pi/p)$.

Thus, based on both cases 1 and 2 $  |\alpha|^2 + |\beta|^2=1,$
  with $\beta = -i\sin(2\pi/p) \neq 0$ since $\sin(2\pi/p) \neq 0,$ for $p \geq 3$.  Thus $t^* = 2\pi/p$ is a valid revival time. Thus the relationship~\eqref{eqq1Thm2revivaltime} holds. 
\end{proof}

\end{document}
